\chapter{Nosebleed (Epistaxis)}

\section*{Definition}
Bleeding from the nasal cavity, arising from the anterior septum (Kiesselbach plexus, 90%) or posterior nasal cavity (more severe, especially in elderly).

\section*{Pathophysiology}
Nasal mucosal vessels are superficial and fragile; anterior epistaxis from Little's area (Kiesselbach plexus) where several arteries anastomose. Posterior epistaxis arises from Woodruff plexus in the inferior meatus. Bleeding may be spontaneous or triggered by trauma, dryness, inflammation, or systemic coagulopathy.

\section*{Etiology}
\begin{itemize}
  \item Nasal trauma or picking
  \item Dry air (low humidity, heating, air conditioning)
  \item Allergic rhinitis or sinusitis
  \item Anticoagulant medications (warfarin, aspirin, clopidogrel)
  \item Nasal tumors or polyps
  \item Coagulation disorders (hemophilia, von Willebrand disease)
  \item Hypertension (especially posterior bleeds)
  \item Pregnancy (hormonal vascular changes)
  \item Cocaine use or intranasal medications
  \item Hereditary hemorrhagic telangiectasia (Osler-Weber-Rendu)
\end{itemize}

\section*{Indications for Evaluation}
Seek care if:
\begin{itemize}
  \item Severe or worsening symptoms
  \item Bleeding persisting more than 20 minutes despite direct pressure, recurrent episodes, heavy blood loss, or bleeding on anticoagulation therapy
  \item Accompanied by other concerning signs
\end{itemize}

\section*{Related Symptoms}
\begin{itemize}
  \item Nasal congestion
  \item Headache
  \item Dizziness
  \item Pale skin (anemia)
  \item Easy bruising
\end{itemize}
\textbf{Note:} Always consider serious underlying conditions when evaluating persistent nosebleed (epistaxis).

\section*{Self-Care / Home Management}
\begin{itemize}
  \item Rest and avoid activities that may aggravate the condition.
  \item Maintain adequate hydration and a balanced diet.
  \item Over-the-counter remedies may provide symptomatic relief (consult a pharmacist or doctor).
  \item Monitor symptoms and seek professional advice if they persist or worsen.
  \item Avoid self-medicating with prescription medications without medical guidance.
\end{itemize}

\section*{Diagnostic Approach}
\begin{itemize}
  \item A thorough history and physical examination are the first steps in evaluation.
  \item Laboratory tests (blood work, urinalysis) may be ordered based on clinical suspicion.
  \item Imaging studies (X-ray, ultrasound, CT, MRI) may be indicated when structural pathology is suspected.
  \item Specialist referral is arranged when the diagnosis remains uncertain or requires specific expertise.
\end{itemize}

\section*{Treatment / Management Overview}
\begin{itemize}
  \item Treatment targets the underlying cause identified through diagnostic evaluation.
  \item Supportive care and symptom management are provided while awaiting definitive therapy.
  \item Medications, lifestyle modifications, and physical therapies may be prescribed as appropriate.
  \item Follow-up ensures treatment effectiveness and allows timely adjustments to the care plan.
\end{itemize}

\clearpage